% Werte und Aufgabenstellung: aufgaben/Blatt07.pdf, S. 2.
\subsection*{Gegeben und Modell}
Für Wasser ist $\sigma=\qty{0.0721}{\newton\per\metre}$ gegeben, für
Kupfer $\rho_\mathrm{Cu}=\qty{8.96}{\gram\per\centi\metre\cubed}
=\qty{8960}{\kilogram\per\metre\cubed}$. Als zusätzliche Näherung
verwenden wir $g=\qty{9.81}{\metre\per\second\squared}$.

Für b und c betrachten wir eine homogene, horizontale Kreisscheibe mit
Radius $R$ und Dicke $d$. Wir interpretieren den Hinweis des Blattes so, dass
an der gesamten kreisförmigen Kontaktlinie eine einzige Flüssigkeitsoberfläche
mit voller Stärke $\sigma$ senkrecht nach oben zieht. Dies ist eine idealisierte
maximale Vertikalkomponente, nicht eine Kraft entlang der Umfangsrichtung.
Hydrostatischen Auftrieb vernachlässigen wir ausdrücklich, um das einfache
Oberflächenkraftmodell des Hinweises zu erhalten. Das ist eine zusätzliche
Modellannahme und keine allgemeine Folge des Hinweises. Bei Kupfer und Wasser
ist das Gewicht der Flüssigkeit mit dem Volumen der Münze nur etwa ein Neuntel
des Münzgewichts; die verdrängte Flüssigkeit und der Meniskus können jedoch
mehr zur Druckkraft beitragen. Ohne Meniskus- und Benetzungsmodell ist die
Vernachlässigung nicht genau überprüfbar. Reale Schwimmgrenzen können daher
abweichen.

\subsection*{a) Druck im kugelförmigen Tropfen}
Ein Tropfen mit Radius $r$ hat eine Flüssigkeitsoberfläche. Für eine kleine
Radiusänderung gilt
\begin{align}
  A&=4\pi r^2, & \mathrm dA&=8\pi r\,\mathrm dr,\\
  V&=\frac43\pi r^3, & \mathrm dV&=4\pi r^2\,\mathrm dr.
\end{align}
Die Netto-Druckarbeit $\Delta p\,\mathrm dV$ gleicht bei quasistatischer
Änderung die Oberflächenenergie $\sigma\,\mathrm dA$ aus. Für $r>0$ folgt
\begin{ergebnis}[a) Überdruck]
\begin{equation*}
  \Delta p=p_\mathrm{innen}-p_\mathrm{aussen}
  =\sigma\frac{\mathrm dA}{\mathrm dV}=\frac{2\sigma}{r},
  \qquad p_\mathrm{innen}=p_\mathrm{aussen}+\frac{2\sigma}{r}.
\end{equation*}
\end{ergebnis}
Das $p$ in der auf dem Blatt angegebenen Arbeitsgleichung ist also der
Überdruck, nicht der absolute Innendruck. Der Aussendruck ist nicht angegeben.
Eine Seifenblase hat dagegen zwei Grenzflächen und deshalb einen anderen Faktor.

\subsection*{b) Zulässige Münzabmessungen}
Das Gewicht und die nach oben gerichtete ideale Oberflächenkraft sind
\begin{equation}
  F_g=\rho_\mathrm{Cu}\pi R^2 d g,
  \qquad F_\sigma=\sigma\,2\pi R.
\end{equation}
Für $R,d>0$ muss $F_g\leq F_\sigma$ gelten. Division durch $\pi R$ gibt
\begin{ergebnis}[b) Bedingung für Radius und Dicke]
\begin{equation*}
  R d\leq\frac{2\sigma}{\rho_\mathrm{Cu}g}
  \approx\qty{1.6405e-6}{\metre\squared}.
\end{equation*}
\end{ergebnis}
Die Grenzwerte liegen auf der Gleichheitskurve. Ein einzelnes unabhängiges
Paar aus maximalem Radius und maximaler Dicke existiert nicht: Wird eines
vorgegeben, erhält man den Grenzwert des anderen,
\begin{equation}
  R_\mathrm{max}(d)=\frac{2\sigma}{\rho_\mathrm{Cu}g d},
  \qquad d_\mathrm{max}(R)=\frac{2\sigma}{\rho_\mathrm{Cu}g R}.
\end{equation}
Dünnere Scheiben können in diesem Modell einen grösseren Radius haben.

\subsection*{c) Fünfräppler}
Für den vorgegebenen Durchmesser $D=\qty{17.15}{\milli\metre}$ gilt
$R=\qty{8.575}{\milli\metre}$; die Dicke ist
$d=\qty{1.25}{\milli\metre}$. Mit der vorgegebenen Kupferdichte ist die
Masse im Scheibenmodell
\begin{equation}
  m=\rho_\mathrm{Cu}\pi R^2d\approx\qty{2.587}{\gram},
  \qquad F_g=mg\approx\qty{0.02538}{\newton}.
\end{equation}
Im Grenzfall $2\pi R\sigma_\mathrm{min}=F_g$ ergibt sich
\begin{ergebnis}[c) Erforderliche Oberflächenspannung]
\begin{equation*}
  \sigma_\mathrm{min}=\frac{\rho_\mathrm{Cu}gRd}{2}
  \approx\qty{0.4711}{\newton\per\metre}.
\end{equation*}
\end{ergebnis}
Für Wasser liefert das Modell nur $F_\sigma=2\pi R\sigma
\approx\qty{0.003885}{\newton}$ und damit rund \qty{15.3}{\percent} des
Münzgewichts. Die Münze schwimmt nach diesem vereinfachten Modell nicht auf
Wasser. Alternativ ergibt b bei dieser Dicke $R_\mathrm{max}\approx
\qty{1.312}{\milli\metre}$, also $D_\mathrm{max}\approx
\qty{2.625}{\milli\metre}$, deutlich weniger als der vorgegebene Durchmesser.

\begin{kontrolle}[Einheiten und Grenzen]
$\sigma/r$ hat die Einheit Druck, $Rd$ die Einheit Fläche und
$\rho gRd$ die Einheit \si{\newton\per\metre}. Der in b und c berechnete
Grenzwert gilt nur bei voller vertikaler Oberflächenkraft und ohne Auftrieb.
Die Bewegung oder Stabilität einer realen Münze kann damit nicht vollständig
vorhergesagt werden.
\end{kontrolle}

\begin{merksatz}
Für Grenzflächenkräfte ist die Kontaktlänge entscheidend, für das Gewicht das
Volumen. Bei Druckarbeit zählt die Druckdifferenz über der Grenzfläche.
\end{merksatz}
